% ID: SEL-D0017-BC-E06
% ESTADO: validada. S018. Revisión matemática, didáctica y visual.
\textbf{a) Probabilidad de haber viajado.}\par\nopagebreak[4]
Hay $15000-8500=6500$ hombres. Sea $V$ haber viajado al extranjero. Los recuentos de personas que viajaron son $8500(0.15)=1275$ mujeres y $6500(0.20)=1300$ hombres. La probabilidad total es
\[
\boxed{P(V)=\frac{1275+1300}{15000}=\frac{103}{600}\approx0.1717.}
\]

\textbf{b) Hombre entre quienes no viajaron.}\par\nopagebreak[4]
No viajaron $12425$ personas, de las cuales $6500-1300=5200$ son hombres. Si $H$ denota ser hombre,
\[
\boxed{P(H\mid V^c)=\frac{5200}{12425}=\frac{208}{497}\approx0.4185.}
\]
El denominador es el grupo que cumple la condición.
